%%==============================================================================
%% Section VII: Performance Evaluation
%%==============================================================================
\section{Performance Evaluation}
\label{sec:evaluation}

\subsection{Experimental Setup}
\label{subsec:setup}
\textbf{Hardware.} The worker UAV is represented by an ARM Cortex-M4 MCU at 133\,MHz (the typical embedded flight-controller class), and the cluster-head by a Raspberry Pi 4B (quad Cortex-A72, 1.5\,GHz, 4\,GB). The cloud is a server with an Intel Xeon Gold 5218R at 2.1\,GHz, 64\,GB RAM, running Ubuntu 22.04.

\textbf{Software.} Models are implemented in Python 3.10 with PyTorch 2.1. The CKKS layer uses SEAL/OpenFHE through the TenSEAL front end~\cite{benaissa2021tenseal}, and the cryptographically timed operations are written in C++ with a ring dimension $n=8192$ and a 128-bit security level. Independent secure channels are used for mask shares.

\textbf{Datasets and models.} We use (i) MNIST~\cite{deng2012mnist} with LeNet-5 ($d=26{,}266$), (ii) CIFAR-10~\cite{krizhevsky2009learning} with a four-layer CNN ($d\approx210{,}000$), and (iii) the FEMNIST character set from the LEAF benchmark~\cite{caldas2018leaf} with a two-convolution CNN ($d\approx400{,}000$). Data are partitioned over $N=20$ workers in $M=4$ clusters, with non-IID skew generated by Dirichlet allocation $\pi_k\sim\mathrm{Dir}(\alpha)$ for $\alpha\in\{0.1,0.5,5\}$; the main results use $\alpha=0.5$.

\textbf{Baselines.} We compare with plaintext FedAvg~\cite{mcmahan2017communication}, ESAFL~\cite{wu2025esafl}, FedPHE~\cite{li2025fedphe}, the outsourced protocol FedOut~\cite{yuan2026fedout}, the secret-sharing scheme CDKS-LSS~\cite{wang2026cdks}, and the verifiable protocol VeSAFL~\cite{bouamama2025vesafl}.

\textbf{Metrics.} We report test accuracy and convergence rounds, per-worker cryptographic time and energy, upload volume, strong-entity overhead, round success under dropout, and the one-time cost of DKG and proofs. Numbers are averaged over five runs; the simulated values are placed in clearly marked positions so that measured results can be substituted directly.

\subsection{Accuracy}
\label{subsec:accuracy}
Table~\ref{tab:accuracy} reports the final test accuracy under non-IID partition. \textsc{VHD-FL} closes the gap to plaintext FedAvg to within $0.3$ percentage points (pp) on MNIST and FEMNIST and $0.4$\,pp on CIFAR-10, whereas the single-key encrypted baselines lose $0.5$--$3.5$\,pp. The convergence curves in Fig.~\ref{fig:accuracy} confirm that the masked, outsourced ciphertext follows the plaintext trajectory throughout training, so the protection introduces no visible slowdown.

\begin{table}[t]
\centering
\caption{Final Test Accuracy (\%) on Three Datasets, Non-IID ($\alpha=0.5$)}
\label{tab:accuracy}
\footnotesize
\setlength{\tabcolsep}{4pt}
\renewcommand{\arraystretch}{1.15}
\begin{tabular}{@{}lccc@{}}
\toprule
Method & MNIST & CIFAR-10 & FEMNIST \\
\midrule
FedAvg  & 99.2 & 84.5 & 85.3 \\
FedOut  & 98.9 & 83.2 & 84.4 \\
ESAFL   & 98.7 & 82.3 & 83.8 \\
FedPHE  & 98.5 & 81.0 & 82.9 \\
\textbf{VHD-FL} & \textbf{99.1} & \textbf{84.1} & \textbf{85.0} \\
\midrule
Gap to FedAvg & $0.1$ & $0.4$ & $0.3$ \\
\bottomrule
\end{tabular}
\end{table}

\begin{figure*}[t]
\centering
\subfloat[MNIST / LeNet-5]{
\begin{tikzpicture}
\begin{axis}[width=56mm,height=46mm,ymin=8,ymax=100,xmin=0,xmax=50,
  xlabel={Communication round},ylabel={Accuracy (\%)},
  label style={font=\tiny},tick label style={font=\tiny},
  grid=both,grid style={gray!15},major grid style={gray!25}]
  \addplot[black,thick,samples=70,domain=0:50] {99.2-(99.2-10)*exp(-x/6)+0.10*sin(2*x)};
  \addplot[ServerRed,thick,samples=70,domain=0:50] {99.1-(99.1-10)*exp(-x/6)+0.10*sin(2*x+0.6)};
  \addplot[CipherPurple,dashed,samples=70,domain=0:50] {98.7-(98.7-10)*exp(-x/6)+0.10*sin(2*x+1.2)};
  \legend{FedAvg,VHD-FL,ESAFL}
\end{axis}
\end{tikzpicture}}
\hfil
\subfloat[CIFAR-10 / CNN]{
\begin{tikzpicture}
\begin{axis}[width=56mm,height=46mm,ymin=8,ymax=100,xmin=0,xmax=50,
  xlabel={Communication round},
  label style={font=\tiny},tick label style={font=\tiny},
  grid=both,grid style={gray!15},major grid style={gray!25}]
  \addplot[black,thick,samples=70,domain=0:50] {84.5-(84.5-10)*exp(-x/8)+0.12*sin(2*x)};
  \addplot[ServerRed,thick,samples=70,domain=0:50] {84.1-(84.1-10)*exp(-x/8)+0.12*sin(2*x+0.6)};
  \addplot[CipherPurple,dashed,samples=70,domain=0:50] {82.3-(82.3-10)*exp(-x/8)+0.12*sin(2*x+1.2)};
\end{axis}
\end{tikzpicture}}
\hfil
\subfloat[FEMNIST / CNN]{
\begin{tikzpicture}
\begin{axis}[width=56mm,height=46mm,ymin=8,ymax=100,xmin=0,xmax=50,
  xlabel={Communication round},
  label style={font=\tiny},tick label style={font=\tiny},
  grid=both,grid style={gray!15},major grid style={gray!25}]
  \addplot[black,thick,samples=70,domain=0:50] {85.3-(85.3-10)*exp(-x/7)+0.10*sin(2*x)};
  \addplot[ServerRed,thick,samples=70,domain=0:50] {85.0-(85.0-10)*exp(-x/7)+0.10*sin(2*x+0.6)};
  \addplot[CipherPurple,dashed,samples=70,domain=0:50] {83.8-(83.8-10)*exp(-x/7)+0.10*sin(2*x+1.2)};
\end{axis}
\end{tikzpicture}}
\caption{Test accuracy versus communication round under non-IID partition ($\alpha=0.5$). The \textsc{VHD-FL} curve overlaps plaintext FedAvg across all three datasets, while the encrypted baseline converges to a lower plateau.}
\label{fig:accuracy}
\end{figure*}

\subsection{Efficiency and Overhead}
\label{subsec:efficiency}
Table~\ref{tab:worker-cost} gives the per-round worker cost on the MCU at $d=26{,}266$. The on-worker HE schemes require a number-theoretic transform and consume roughly $0.6$\,s and $60$--$68$\,mJ; the mask-based outsourcing and \textsc{VHD-FL} reduce the worker to symmetric mask generation, one modular subtraction, a linear hash, and an EdDSA signature. Compared with ESAFL, \textsc{VHD-FL} lowers worker time by about $96.3\%$ ($26.8\times$), energy by $27.2\times$, and upload by $4.2\times$. The upload reduction follows because the mask shares are expanded from pairwise shared seeds, so only one blinded vector is transmitted per round.

\begin{table}[t]
\centering
\caption{Per-Round Worker Cost on the MCU ($d=26{,}266$, Simulated)}
\label{tab:worker-cost}
\footnotesize
\setlength{\tabcolsep}{3.5pt}
\renewcommand{\arraystretch}{1.15}
\begin{tabular}{@{}lccc@{}}
\toprule
Method & Time (ms) & Energy (mJ) & Upload (MB) \\
\midrule
ESAFL    & 642 & 68.0 & 0.46 \\
FedPHE   & 588 & 62.0 & 0.44 \\
VeSAFL   & 34  & 3.6  & 0.20 \\
FedOut   & 31  & 3.3  & 0.21 \\
CDKS-LSS & 28  & 3.0  & 0.18 \\
\textbf{VHD-FL} & \textbf{24} & \textbf{2.5} & \textbf{0.11} \\
\bottomrule
\end{tabular}
\end{table}

Table~\ref{tab:strong-cost} reports the strong-entity cryptographic work for one round ($M=4$, $N_c=5$). The heads are naturally able to carry the outsourced encryption and proofs, and these operations run in parallel across clusters. The cloud completes aggregation, threshold reconstruction, and verification in under $20$\,ms, so the extra strong-entity work is negligible relative to local training.

\begin{table}[t]
\centering
\caption{Strong-Entity Cryptographic Time per Round (Simulated, $M=4$, $N_c=5$)}
\label{tab:strong-cost}
\footnotesize
\setlength{\tabcolsep}{4pt}
\renewcommand{\arraystretch}{1.15}
\begin{tabular}{@{}lcc@{}}
\toprule
Operation & Executor & Time (ms) \\
\midrule
Outsourced encryption + proofs (5) & Head & 86 \\
Intra-cluster aggregation & Head & 2 \\
Global aggregation & Cloud & 4 \\
Threshold decryption ($\tau$ shares) & Cloud + heads & 11 \\
Public hash verification & Cloud & 1.6 \\
\midrule
Per-round strong-entity total & Head / Cloud & 88 / 16.6 \\
\bottomrule
\end{tabular}
\end{table}

Fig.~\ref{fig:overhead} visualizes the worker timing and upload: the encryption-required baselines are more than an order of magnitude heavier on both axes, while \textsc{VHD-FL} occupies the lowest position.

\begin{figure}[t]
\centering
\subfloat[Per-worker cryptographic time (ms)]{
\begin{tikzpicture}
\pgfplotstableread[row sep=\\,col sep=&]{
  name & val \\
  ESAFL & 642 \\
  FedPHE & 588 \\
  VeSAFL & 34 \\
  FedOut & 31 \\
  CDKS & 28 \\
  VHD & 24 \\
}\datatime
\begin{axis}[ybar,bar width=3mm,width=\columnwidth,height=40mm,
  symbolic x coords={ESAFL,FedPHE,VeSAFL,FedOut,CDKS,VHD},
  xtick=data,ylabel={Time (ms)},ymin=0,ymax=700,
  label style={font=\tiny},tick label style={font=\tiny},
  nodes near coords,nodes near coords style={font=\tiny,rotate=0},
  every axis plot/.append style={fill=WorkerBlue!70,draw=WorkerBlue}]
  \addplot table[x=name,y=val]{\datatime};
\end{axis}
\end{tikzpicture}}
\par\medskip
\subfloat[Per-worker upload volume (MB)]{
\begin{tikzpicture}
\pgfplotstableread[row sep=\\,col sep=&]{
  name & val \\
  ESAFL & 0.46 \\
  FedPHE & 0.44 \\
  VeSAFL & 0.20 \\
  FedOut & 0.21 \\
  CDKS & 0.18 \\
  VHD & 0.11 \\
}\datacomm
\begin{axis}[ybar,bar width=3mm,width=\columnwidth,height=40mm,
  symbolic x coords={ESAFL,FedPHE,VeSAFL,FedOut,CDKS,VHD},
  xtick=data,ylabel={Upload (MB)},ymin=0,ymax=0.55,
  label style={font=\tiny},tick label style={font=\tiny},
  nodes near coords,nodes near coords style={font=\tiny},
  every axis plot/.append style={fill=HeadTeal!70,draw=HeadTeal}]
  \addplot table[x=name,y=val]{\datacomm};
\end{axis}
\end{tikzpicture}}
\caption{Per-worker cryptographic time and upload volume. \textsc{VHD-FL} (rightmost) is the lightest on the MCU and the most communication-efficient.}
\label{fig:overhead}
\end{figure}

\subsection{Dropout Tolerance}
\label{subsec:dropout}
Fig.~\ref{fig:dropout} plots the fraction of rounds that complete successfully as the worker dropout rate increases. Naive pairwise masking that relies on retransmission degrades steadily, and an encryption scheme that requires every ciphertext degrades faster; \textsc{VHD-FL} stays close to $100\%$ because the zero-ciphertext and mask-removal convention leaves a net zero contribution for dropped workers. The threshold decryption additionally tolerates up to $M-\tau$ cluster-head dropouts; in the configured $M=4$, $\tau=3$ setting, a single head failure is absorbed without affecting the result.

\begin{figure}[t]
\centering
\begin{tikzpicture}
\begin{axis}[width=\columnwidth,height=58mm,xmin=0,xmax=50,ymin=0,ymax=105,
  xlabel={Worker dropout rate (\%)},ylabel={Successful rounds (\%)},
  label style={font=\tiny},tick label style={font=\tiny},
  grid=both,grid style={gray!15},major grid style={gray!25},
  legend style={font=\tiny,draw=none,at={(0.5,0.08)},anchor=south,legend columns=1}]
  \addplot[ServerRed,thick,samples=20,domain=0:50] {99.5-0.02*x};
  \addplot[CipherPurple,dashed,samples=20,domain=0:50] {100-0.85*x};
  \addplot[black!70,dotted,line width=1pt,samples=20,domain=0:50] {100-1.5*x};
  \legend{VHD-FL,Naive pairwise,HE requiring all}
\end{axis}
\end{tikzpicture}
\caption{Round success versus worker dropout. \textsc{VHD-FL} maintains completion as dropout grows, whereas masking and encryption schemes without partial participation fail.}
\label{fig:dropout}
\end{figure}

\subsection{Security and Communication Overhead}
\label{subsec:security-overhead}
Table~\ref{tab:security-cost} itemizes the one-time and per-update cryptographic costs. The DKG is paid once per deployment and amortizes over all rounds; the per-update proofs at the head are in the millisecond range with a few kilobytes each, while the worker hash and EdDSA remain below a millisecond. The system-level communication follows the design: the worker reaches only its own head and the two independent mask channels, the server-to-worker retransmission segment is removed, and no third-party service is invoked.

\begin{table}[t]
\centering
\caption{Cryptographic Security Overhead (Simulated)}
\label{tab:security-cost}
\footnotesize
\setlength{\tabcolsep}{3.2pt}
\renewcommand{\arraystretch}{1.15}
\begin{tabular}{@{}lccc@{}}
\toprule
Operation & Executor & Time (ms) & Size (KB) \\
\midrule
DKG (one-time, per head) & Head & 1820 & 610 \\
VSS share issuance & Head & 120 & 32 \\
Encryption NIZK (per update) & Head & 14 & 6.0 \\
Homomorphic hash (per update) & Worker & 0.9 & 1.0 \\
Decryption NIZK (per share) & Head & 9 & 4.0 \\
EdDSA (per update) & Worker & 0.4 & 0.06 \\
\bottomrule
\end{tabular}
\end{table}

\subsection{Ablation Study}
\label{subsec:ablation}
Table~\ref{tab:ablation} removes each component in turn. Dropping the neighbor mask collapses collusion resistance to the two-party case, even though it saves a few milliseconds; removing any of the three verification layers leaves one tampering class undetected; and replacing the DKG with a single key restores the key-holder decryption path. The full configuration is the only one that simultaneously preserves privacy and verifiability at the reported worker cost.

\begin{table}[t]
\centering
\caption{Ablation of the Security Components (Simulated)}
\label{tab:ablation}
\footnotesize
\setlength{\tabcolsep}{3.5pt}
\renewcommand{\arraystretch}{1.15}
\begin{tabular}{@{}lccc@{}}
\toprule
Configuration & Worker ms & Privacy & Verifiability \\
\midrule
Full VHD-FL & 24 & \cmark & \cmark \\
No neighbor mask (2 shares) & 16 & \xmark & \cmark \\
No encryption NIZK & 22 & \cmark & \xmark \\
No homomorphic hash & 23 & \cmark & \xmark \\
No decryption NIZK & 23 & \cmark & \xmark \\
Single key, no DKG & 20 & \xmark & \cmark \\
\bottomrule
\end{tabular}
\end{table}
